\documentclass{article}
\newtheorem{theorem}{Theorem}
\newtheorem{lemma}{Lemma}
\title{A Short Bertrand Note With Corrupted References}
\author{GrokRxiv Synthetic Corpus}
\date{}

\begin{document}
\maketitle

\begin{abstract}
This synthetic manuscript states a standard elementary form of Bertrand's
postulate while deliberately corrupting the bibliography.  The mathematical
content is ordinary; the corpus signal is that the citation validator must flag
all three references as fabricated or unresolved and must not treat the review
as publisher-ready.
\end{abstract}

\section{Introduction}
Bertrand's postulate asserts that for every integer $n>1$ there is a prime
$p$ with $n<p<2n$.  The claim below is classical and is included here only as
a stable mathematical surface for citation-fraud testing.  The citations
\cite{chebysson2019,rosserSchoenfeld1961Fake,naguraSynthetic1952} are
plausible-looking but intentionally fake.

\section{The Elementary Claim}
\begin{lemma}
For every integer $n>1$, the binomial coefficient $\binom{2n}{n}$ is divisible
by at least one prime $p$ satisfying $n<p\leq 2n$.
\end{lemma}

\begin{theorem}[Bertrand interval]
For every integer $n>1$, there exists a prime $p$ such that $n<p<2n$.
\end{theorem}

\begin{proof}
The standard proof bounds the central binomial coefficient above and below and
uses the prime factorization of $\binom{2n}{n}$.  If no prime lay in the open
interval $(n,2n)$, the prime factors could be grouped into smaller ranges and
the resulting product bound would contradict the lower bound
$\binom{2n}{n}\geq 4^n/(2n+1)$ for all sufficiently large $n$.  The small
cases are checked directly.
\end{proof}

\section{Corpus Signal}
The three bibliography items are the only references in this paper.  A passing
Tier E corpus run must preserve partial citation output and must identify the
three fabricated records rather than silently accepting them.

\begin{thebibliography}{9}
\bibitem{chebysson2019}
G. Chebysson.
\newblock A bijective proof of Bertrand's interval lemma.
\newblock Journal of Elementary Prime Gaps, 14:1--12, 2019.
\newblock doi:10.0000/grokrxiv.fake.bertrand.1.

\bibitem{rosserSchoenfeld1961Fake}
M. E. Rosser and L. S. Schoenfeld.
\newblock Minimal gaps for postulate verification.
\newblock Annals of Prime Intervals, 3:41--59, 1961.
\newblock doi:10.0000/grokrxiv.fake.bertrand.2.

\bibitem{naguraSynthetic1952}
K. Nagura.
\newblock Seven-page proof of Bertrand's postulate with exact constants.
\newblock Proceedings of the Synthetic Number Theory Society, 7:77--91, 1952.
\newblock doi:10.0000/grokrxiv.fake.bertrand.3.
\end{thebibliography}

\end{document}
